%--------------------------------------------------------------------------------
%
%
%--------------------------------------------------------------------------------
\pagebreak
\section*{CHECK - test a condition}
\addcontentsline{toc}{section}{CHECK - test a condition}

%--------------------------------------------------------------------------------
\begin{iPageDescEntry}{Syntax}
\texttt{check  <expr>  [  , <msgStr>  ]}
\end{iPageDescEntry}

%--------------------------------------------------------------------------------
\begin{iPageDescEntry}{Description}

The \texttt{CHECK} command evaluates the specified expression and reports whether the condition passed or failed. Unlike the \texttt{ASSERT} command, a failed check does not generate an assertion failure or interrupt execution. Instead, the result is simply reported to the user.

If the expression evaluates to true, the command reports a \texttt{PASS} condition. If the expression evaluates to false, the command reports a \texttt{FAIL} condition. An optional message string may be supplied to describe the purpose of the check or identify the condition being tested. If the message text is omitted, the text is taken from the environment variable \texttt{ENV\_CHECK\_DEF\_MSG}.

The \texttt{CHECK} command is useful for lightweight testing, validation, and interactive diagnostics where failures should be reported without stopping execution.

\end{iPageDescEntry}

%--------------------------------------------------------------------------------
\begin{iPageDescEntry}{Example}

The following example evaluates two expressions. The first condition passes, while the second condition fails and displays the associated message string.

\begin{lstlisting}[style=iPageOpStyle]
(4) -> check ( 2 < 3 )
PASS
(5) -> check ( 3 < 2 ), "A simple test"
"A simple test" : FAIL
(6) ->
\end{lstlisting}
\end{iPageDescEntry}

%--------------------------------------------------------------------------------
\begin{iPageDescEntry}{Notes}

\texttt{CHECK} reports failures but does not stop execution. The optional message string can be used to identify individual tests. The expression is considered successful if it evaluates to a non-zero or true value.

// ??? how about an ENV variable with a default message text ?

\end{iPageDescEntry}